\documentclass[10pt,letterpaper]{article}
\usepackage{cornell-notes}

\title{Clause 17.09.03: Class Exception}
\author{}
\date{}
\setDocTitle{Clause 17.09.03: Class Exception}
\setDocSubtitle{C++ 2024}
\setDocOwner{C++ 2024 Cornell Notes}
\setDocDate{\today}
\setCornellCollection{C++ 2024 Cornell Notes}
\setCornellUnitType{Clause}
\setCornellUnitNumber{17.09.03}
\setCornellUnitTitle{Class Exception}
\hypersetup{pdftitle={Clause 17.09.03: Class Exception},pdfsubject={Cornell notes},pdfauthor={}}

\begin{document}
\maketitle
\begin{CornellOverviewBox}{Overview}
\textbf{Classification:} Standard library - Normative\\[2pt]
\textbf{Learning objective:} Understand the rules, terminology, and semantic consequences associated with Class exception.
\end{CornellOverviewBox}\begin{CornellSummaryBox}{Summary}
{\sffamily\bfseries\color{Accent} Summary}\par\vspace{3pt}Section 17.9.3 focuses on Class exception. Apply the specified interface together with its mandates, effects, result conditions, exception behavior, and complexity constraints. When applying it, preserve the distinction between mandatory requirements, recommendations, permissions, and behavior deliberately left undefined, unspecified, or implementation-defined.
\end{CornellSummaryBox}

\end{document}
